% Shared preamble for the launch monitor technology review
\usepackage[T1]{fontenc}
\usepackage[utf8]{inputenc}
\usepackage{lmodern}
\usepackage{microtype}
\usepackage[margin=1in]{geometry}
\usepackage{amsmath,amssymb,bm}
\usepackage{siunitx}
\usepackage{graphicx}
\usepackage{booktabs}
\usepackage{longtable}
\usepackage{array}
\usepackage{multirow}
\usepackage{caption}
\usepackage{subcaption}
\usepackage{enumitem}
\usepackage{xcolor}
\usepackage{tikz}
\usetikzlibrary{arrows.meta,positioning,calc,angles,quotes}
\usepackage{fancyhdr}
\usepackage{titlesec}
\usepackage[hidelinks,pdfusetitle]{hyperref}
\usepackage[capitalise,noabbrev]{cleveref}

% Bibliography: structured BibTeX database (references.bib), biber backend.
% Entries are grouped in the printed bibliography by their `keywords`
% field -- see CONVENTIONS.md before adding references.
\usepackage[backend=biber,style=numeric,sorting=nyt,giveninits=true,
            maxbibnames=6,minbibnames=6,isbn=false,eprint=false]{biblatex}
\addbibresource{references.bib}
% Print the URL note field compactly for online/manual sources.
\setlength{\bibitemsep}{0.5ex}

% Color scheme
\definecolor{accentblue}{RGB}{20,60,110}
\definecolor{ofgray}{RGB}{90,95,100}
\definecolor{accentgreen}{RGB}{25,110,60}

\titleformat{\chapter}[display]
  {\normalfont\huge\bfseries\color{accentblue}}
  {\chaptertitlename\ \thechapter}{12pt}{\Huge}
\titlespacing*{\chapter}{0pt}{0pt}{24pt}
\titleformat*{\section}{\Large\bfseries\color{accentblue}}
\titleformat*{\subsection}{\large\bfseries\color{ofgray}}

\pagestyle{fancy}
\fancyhf{}
\fancyhead[L]{\small\itshape Launch Monitor Technology}
\fancyhead[R]{\small\itshape\nouppercase{\leftmark}}
\fancyfoot[C]{\thepage}
\renewcommand{\headrulewidth}{0.4pt}

\sisetup{per-mode=symbol,range-phrase=--,range-units=single}

% Convenience macros
\newcommand{\degs}{\si{\degree}}
\newcommand{\mph}{\,mph}
\newcommand{\rpm}{\,rpm}
\newcommand{\patent}[1]{\href{https://patents.google.com/patent/#1}{#1}}
\newcommand{\vect}[1]{\bm{#1}}
\newcommand{\uvec}[1]{\hat{\bm{#1}}}

% Callout box for design implications. Generic by design: this box flags a
% consequence an implementer must act on, independent of any particular product.
\usepackage[most]{tcolorbox}
\newtcolorbox{implication}{
  colback=accentgreen!6, colframe=accentgreen!70!black,
  title=Design implication, fonttitle=\bfseries,
  boxrule=0.6pt, arc=2pt, left=6pt, right=6pt, top=4pt, bottom=4pt,
  breakable
}
\newtcolorbox{keypoint}{
  colback=accentblue!5, colframe=accentblue!70,
  title=Key point, fonttitle=\bfseries,
  boxrule=0.6pt, arc=2pt, left=6pt, right=6pt, top=4pt, bottom=4pt,
  breakable
}
% Callout for widely-repeated claims that are wrong, contested, or unsourced.
% Used to flag corrections so they are not silently absorbed on later edits.
\newtcolorbox{warning}{
  colback=red!4, colframe=red!55!black,
  title=Caution, fonttitle=\bfseries,
  boxrule=0.6pt, arc=2pt, left=6pt, right=6pt, top=4pt, bottom=4pt,
  breakable
}
